% ---------------------------------------------------------------------------
% Contrastive Learning for Answer-Aware SPARQL Query Embeddings
% ESWC 2027, Research track. Springer LNCS (llncs.cls + splncs04.bst).
% Open items are typeset in red with \hp{...}; search for "\hp{".
% ---------------------------------------------------------------------------
\documentclass[runningheads]{llncs}

\usepackage[T1]{fontenc}
\usepackage[utf8]{inputenc}
\usepackage{graphicx}
\usepackage{booktabs}
\usepackage{hhline}
\usepackage{amsmath,amssymb}
\usepackage{listings}
\usepackage[table]{xcolor}
% table highlight colours
\definecolor{headblue}{RGB}{221,232,245}
\definecolor{grpgray}{RGB}{242,242,242}
\definecolor{sumgray}{RGB}{230,230,230}
\definecolor{bestgreen}{RGB}{204,236,208}
\definecolor{posgreen}{RGB}{220,242,222}
\definecolor{negred}{RGB}{250,222,220}
\definecolor{ftblue}{RGB}{235,242,250}
\usepackage{url}
\usepackage[hidelinks]{hyperref}
\pdfstringdefDisableCommands{\def\\{ }}
\usepackage{tikz}
\usetikzlibrary{arrows.meta,positioning,fit,calc,shapes.geometric,backgrounds}
\usepackage{pgfplots}
\pgfplotsset{compat=1.18}
\usepackage{adjustbox}
\usepackage{placeins}
\lstset{
  basicstyle=\ttfamily\scriptsize,
  breaklines=true,
  frame=single,
  columns=fullflexible,
  keepspaces=true,
  aboveskip=2pt,
  belowskip=2pt,
}

\newcommand{\squareds}{\textsc{SQuARE}}
\newcommand{\Ans}{\mathrm{Ans}}
% grey gain annotation used in the compact tables
\newcommand{\dg}[1]{{\scriptsize\color{black!80}#1}}
% open item (TODO): rendered in red
\newcommand{\hp}[1]{{\color{red!75!black}#1}}

\setlength{\emergencystretch}{2em}
\begin{document}

\title{Contrastive Learning for Answer-Aware\\SPARQL Query Embeddings}

\titlerunning{Answer-Aware SPARQL Query Embeddings}

\author{Anonymous Author(s)}
\authorrunning{Anonymous}
\institute{Anonymous Institution}

\maketitle

% ===========================================================================
\begin{abstract}
Pretrained SPARQL query representations are not explicitly optimized for
answer overlap on a target RDF graph, although query recommendation,
query-log analysis and approximate result reuse need this signal without
executing each incoming query. Syntax alone is an imperfect proxy: similar queries may return disjoint answers, and different queries may
return identical ones. We introduce answer-supervised representation learning
for SPARQL, which exploits the fact that a query's denotation can be observed
automatically. Each query of a log is executed once, offline, and query pairs
are labelled by the Jaccard overlap of their answers; at inference time, only
the query text is needed. Based on the LSQ logs of DBpedia and Wikidata, we
release \squareds{}, a resource with cached answer sets, 92{,}664 positive and
282{,}720 negative query pairs, and splits in which no positive relationship
crosses from training to test. A paired before-and-after study of ten pretrained encoders from four model families shows that answer-supervised fine-tuning improves retrieval among held-out queries for every encoder (on DBpedia, Hits@1 rises from 0.34 to 0.55 on
average) and substantially reduces the spread between encoders. The study
also shows that ranking quality and similarity calibration are distinct
problems: when the embeddings decide approximate result reuse with a fixed
threshold, 59\% to 77\% of the reused answers on DBpedia and at most 28\% on
Wikidata meet the overlap criterion, so a learned similarity alone does not
make reuse reliable. The planned code and data release is \url{https://anonymous.4open.science/r/SQuARE-ESWC27}.

\keywords{SPARQL \and Query Embeddings \and Contrastive Learning \and
Semantic Caching \and Query Logs \and Knowledge Graphs}
\end{abstract}

% ===========================================================================
\section{Introduction}
\label{sec:intro}

SPARQL embeddings map queries to vectors so that we can find related queries
by vector similarity. They promise a fast way to search query logs such as
LSQ~\cite{saleem2015lsq}, which collect millions of queries sent to public
endpoints: given a new query, an embedding retrieves earlier queries without
executing it. Two tasks build on this search. In \emph{query retrieval}, a
user or a system looks for earlier queries that return related data; in
\emph{approximate result reuse}, a system returns a stored answer instead of
executing the new query, as in semantic
caching~\cite{dar1996semanticcaching,papailiou2015caching}. Both tasks
depend on one specific relation: two queries should be close when they
return similar answers on the target graph.
Textual similarity does not necessarily reflect this relation. Figure~\ref{fig:gap}
shows three DBpedia queries. Query (b) differs from (a) in
a single predicate, \texttt{deathPlace} instead of \texttt{birthPlace}, yet
the two queries share only 37 answers ($J = 0.19$). Query (c), by contrast,
adds a variable and a \texttt{FILTER} but returns exactly the same answer set as (a). By token overlap, however, (b) is closer to (a) than (c) is, so a similarity based on text alone ranks (b) first (Figure~\ref{fig:gap}, panel~2). An embedding that serves our tasks should do the opposite and place (c) closer to (a) than (b) (panel~3).
Capturing this ordering requires evidence beyond the query text, because
answer overlap depends on the data of the graph~\cite{perez2009sparql}.

\begin{figure}[t]
\centering
\input{fig-teaser}
\caption{\small{(1)~Query~(b)
changes one predicate of~(a) and shares few answers with it on DBpedia,
whereas query~(c) rewrites~(a) and returns the same answers. By token
similarity, however, (b) is closer to~(a) than~(c) is. (2)~A text-based
encoder places queries by how they look and therefore retrieves~(b) for~(a).
(3)~Our answer-aware encoder, trained on answer overlap observed offline,
places queries by what they return and retrieves~(c), the right match for
query retrieval and approximate result reuse.}}
\label{fig:gap}
\end{figure}

Existing embedding approaches do not use such evidence. Their supervision is
\emph{result-blind}: it reflects how a query is written, not what it returns.
Query2Vec and Querc embed SQL queries for workload tasks such as
summarization and error prediction~\cite{jain2018query2vec,jain2019querc},
SPBERT pretrains a language model on SPARQL logs~\cite{tran2021spbert}, and
general text and code encoders such as E5 and CodeBERT learn from other
data~\cite{wang2022e5,feng2020codebert}. These models may capture useful
regularities in query text, but their training never aligns embedding
similarity with answer overlap.

Other approaches address related but distinct problems. Containment and
equivalence checking establish relations that hold on every graph; these
problems are computationally hard even for small SPARQL
fragments~\cite{pichler2014containment}, and such universal relations miss
the partial, data-dependent overlap between (a) and (b). Similarity search
over RDF graphs~\cite{zheng2016sparqlsimilarity,calisto2025graphsimilarity}
retrieves graph content rather than queries from a log. SQA compares
execution results to assess a generated query against a
reference~\cite{bao2026sqa}, but does not use them to train a query encoder.
None of these lines of work answers the question that our tasks raise:
\emph{can answers collected offline teach an encoder to retrieve
high-overlap queries from text alone?}

We address this question with a simple observation: query answers need no
human annotation, because an endpoint can compute them. This lets us
separate the collection of data-dependent evidence from its use at query
time. We therefore propose \emph{answer-supervised representation learning}
for SPARQL query embeddings. We execute each logged query once, offline, and
label query pairs by the Jaccard overlap of their answer sets: pairs with
high overlap serve as positives, and pairs with little or no overlap serve
as negatives.
We then introduce a contrastive objective that aligns embedding distances
with these labels (Section~\ref{sec:method}). It pulls high-overlap queries
together even when their formulations differ, and it pushes low-overlap
queries apart even when their text is similar, which is exactly the ordering
that Figure~\ref{fig:gap} calls for. Because the labels come from the target
graph, the encoder learns overlap on that graph rather than relations that
hold on every graph. Execution thus supervises training only: at inference,
the encoder reads the query text alone.
We evaluate this approach on \squareds{}, which we build from the DBpedia
and Wikidata logs in LSQ, with ten pretrained encoders from four model
families. Fine-tuning improves retrieval among held-out queries for
every encoder on both graphs, raising mean Hits@1 on DBpedia from $0.34$ to
$0.55$, and it narrows the gap between strong and weak encoders. Result
reuse yields a second finding: better ranking alone does not make reuse
reliable. With a fixed similarity threshold, most reused answers on Wikidata
fall below the overlap criterion, so reliable reuse also needs calibrated
decisions (Section~\ref{sec:caching}). We summarize our contributions as follows:
\begin{itemize}
  \item We define query relatedness by answer
        overlap on a target graph and train SPARQL query embeddings from
        offline execution, with query text as the only input at inference
        (Section~\ref{sec:method}).
  \item We release \squareds{}, which provides cached
        answer sets for 15{,}476 queries, 92{,}664 positive pairs and
        282{,}720 negative pairs from DBpedia and Wikidata, with splits that
        keep every positive relation within one split
        (Section~\ref{sec:data}), together with our code.
  \item We compare ten encoders before and after
        fine-tuning on query retrieval and approximate result reuse
        (Section~\ref{sec:exp}).
\end{itemize}

%The remainder of this paper is organized as follows. Section~\ref{sec:related} reviews related work, Section~\ref{sec:method} presents the method and the resource, Section~\ref{sec:exp} reports the experiments, and Section~\ref{sec:conclusion} concludes.
% ===========================================================================
\section{Related Work}
\label{sec:related}

\paragraph{Contrastive representation learning.}
InfoNCE~\cite{oord2018cpc} and its supervised
variant~\cite{khosla2020supcon} underlie modern sentence and passage
encoders~\cite{reimers2019sbert,gao2021simcse,karpukhin2020dpr,wang2022e5}.
Text encoders derive positives from paraphrase pairs or dropout
noise~\cite{gao2021simcse} and mine negatives from their own
index~\cite{xiong2021ance}; both are approximations, since mined negatives may
be false negatives. SPARQL queries differ in that their answers can be computed.
Positives can therefore be verified by execution, and negatives can be
verified to share almost no answers with the anchor. These labels define the
resource; in-batch candidates from other anchors are not verified for every
anchor, but known positives are masked (Section~\ref{sec:objective}). Our objective is
standard; what is new is the task, graph-relative answer-overlap retrieval for
SPARQL, and supervision for it that is measured rather than annotated.

\paragraph{Learning from execution.}
Using execution as a training signal is established for program code.
CodeScore~\cite{dong2023codescore} learns to predict the functional
correctness of generated code from execution outcomes, and
CodeScore-R~\cite{yang2024codescorer} combines contrastive learning with
semantics-preserving program transformations to make such a metric robust.
Both compare a program with a reference specification on test inputs, whereas
we compare two queries by the answers they return on one fixed graph, a
relation that depends on the data and not only on program behavior.

\paragraph{Encoders for queries and code.}
CodeBERT~\cite{feng2020codebert} and
GraphCodeBERT~\cite{guo2021graphcodebert} are pretrained on source code, the
latter incorporating data-flow structure. SPBERT~\cite{tran2021spbert} is
pretrained on SPARQL, but for question-to-query generation, and is therefore
never exposed to a query-to-query signal. General-purpose text
embedders~\cite{song2020mpnet,wang2022e5,xiao2024bge,li2023gte,zhang2025qwen3emb}
and code language models~\cite{hui2024qwen25coder} are frequently applied to
queries as plain strings. None of these models is trained to predict answer
overlap on a target RDF graph. In this work, we treat them as inputs to our
method rather than as competing approaches.

\paragraph{Query embeddings for logs and workloads.}
Learned representations of whole queries have mainly been studied for
relational workloads. Query2Vec~\cite{jain2018query2vec} and
Querc~\cite{jain2019querc} embed SQL query text with NLP models to support
workload tasks such as error prediction, user labelling and index
recommendation; their training signal is the query text itself. For SPARQL,
embeddings have mostly been applied to the graph rather than to the query:
TrQuery~\cite{zhang2018trquery} embeds the entities and relations of the RDF
graph to score approximate solutions of a single query. Neither line compares
two queries by the answers they return, which is the relation we learn.

\paragraph{SPARQL equivalence, query logs and benchmarks.}
Query containment and equivalence are classical decision
problems~\cite{chandra1977,pichler2014containment}. They are exact but
computationally expensive, and they are defined over all databases, whereas
our notion of relatedness is relative to a specific graph.
LSQ~\cite{saleem2015lsq} and USEWOD~\cite{luczak2014usewod} publish real
query logs with structural metadata, but do not define relatedness between
queries. Performance
benchmarks~\cite{morsey2011dbpsb,saleem2015feasible,angles2022wdbench} select
queries to stress query engines and use runtime as ground truth. Question
answering
benchmarks~\cite{dubey2019lcquad,usbeck2017qald7,kosten2023spider4sparql}
pair a natural-language question with a single gold query and thus support
semantic parsing rather than query-to-query retrieval. To the best of our
knowledge, no existing SPARQL resource provides query-pair supervision
grounded in execution.

\paragraph{Semantic caching.}
Semantic caching~\cite{dar1996semanticcaching} reuses a cached result when a
stored query matches or subsumes an incoming one, a decision that is made
symbolically. For SPARQL, Papailiou et al.~\cite{papailiou2015caching} cache
the results of subgraph patterns and match them structurally against incoming
queries, and Lorey and Naumann~\cite{lorey2013templates} detect query
templates in logs to prefetch data. Both detect reuse that is visible in the
query structure, but not data-dependent overlap. Semantic caches for large
language models, such as GPTCache~\cite{bang2023gptcache}, instead reuse an
answer when the embedding of a new request is close enough to a cached one.
Our reuse experiment follows this design with an embedding learned from
execution; since the decision thresholds the similarity, its absolute
\emph{scale} matters in addition to its ranking.

% ===========================================================================
\section{Methodology}
\label{sec:method}

Our method has three steps (Figure~\ref{fig:pipeline}). We execute each
logged query once and cache its answers, label query pairs by the overlap of
these answers and split them without leakage (Section~\ref{sec:data}), and
fine-tune a pretrained encoder on the labeled pairs with a contrastive
objective (Section~\ref{sec:objective}). We first define the task.

\begin{figure}[t]
\centering
\input{fig-pipeline}
\caption{Overview of our approach. Offline: (1)~execute each logged query
once and cache its answers; (2)~label query pairs by answer overlap $J$;
(3)~split by connected components of positive pairs; (4)~fine-tune an
encoder contrastively. Online: the frozen encoder embeds only the query text
and ranks cached queries for retrieval or approximate reuse.}
\label{fig:pipeline}
\end{figure}

\subsection{Problem Formulation}
\label{sec:problem}

\paragraph{Query fragment and answers.}
We consider SPARQL \texttt{SELECT} queries; \texttt{ASK}, \texttt{CONSTRUCT}
and \texttt{DESCRIBE} queries are out of scope. Let $G$ be an RDF graph. We
denote by $\Ans(q,G)$ the \emph{set} of solution rows returned by $q$ over
$G$. Each row is canonicalized as the tuple of its RDF terms in projection
order, retaining the term type, lexical value, language tag and datatype of
each term, while variable names are discarded, so that variable renaming does
not alter the answer. Duplicate rows are collapsed and row order is ignored.
Without slicing, \texttt{DISTINCT} and \texttt{ORDER BY} therefore do not
affect $\Ans(q,G)$. With \texttt{LIMIT} or \texttt{OFFSET}, however, both
modifiers can change \emph{which} rows are returned, since ordering and
duplicate elimination precede slicing~\cite{w3c2013sparql11}; $\Ans(q,G)$ is
then the set of rows actually returned for the logged query. This case is
rare: no DBpedia query and 88 Wikidata queries (0.8\%) use \texttt{LIMIT} or
\texttt{OFFSET}. Aggregates, which occur in 18\% of the Wikidata queries and in
none of the DBpedia queries, are treated like any other projected value.
Each query was executed with a timeout of 120\,s and up to three retries;
queries that failed, timed out or returned no rows were discarded.
We found no sign of a row limit: in the test split, answer sets reach 2.8
million rows on DBpedia and 5.6 million on Wikidata, and no answer size
recurs as a cap would cause. Blank-node labels are
scoped to a single result document~\cite{w3c2013sparqljson}, so they are not
comparable across queries; we \hp{[TODO: the current canonicalizer keeps
blank-node labels as values. Count affected queries with
table1\_from\_cache.py; if non-zero, replace each blank node by a placeholder
term and relabel, otherwise state that no answer contains a blank node.]}.
Two consequences follow. First, $\Ans$ is a set-based representation:
$J = 1$ means identical canonical answer sets on $G$, not interchangeable
SPARQL responses, since multiplicities, order and variable names are
discarded. Second, returning one query's rows for another additionally
requires mapping its output columns to the other query's variables.

\paragraph{Result similarity.}
The \emph{result similarity} of two queries is the Jaccard overlap of their
answer sets:
\begin{equation}
J(q_i,q_j) = \frac{|\Ans(q_i,G) \cap \Ans(q_j,G)|}{|\Ans(q_i,G) \cup \Ans(q_j,G)|}.
\label{eq:jaccard}
\end{equation}
A pair is \emph{high-overlap} (positive) if $J \geq 0.8$ and
\emph{low-overlap} (negative) if $J \leq 0.1$; intermediate pairs are not
assigned explicit positive or negative labels. Cross-triplet batching can
nevertheless introduce such pairs as negative candidates
(Section~\ref{sec:objective}). High overlap is not an equivalence relation: it is not
transitive, and it admits pairs whose answers differ in a few rows. In \squareds{},
79\% (DBpedia) and 86\% (Wikidata) of the positive pairs have $J = 1$. The distinction matters
little for retrieval, but it is essential when one query's answer is returned
in place of another's (Section~\ref{sec:caching}).
We only retain queries with non-empty answers. $J$ is undefined for two empty
answer sets, and empty answers can stem from errors such as misspelled IRIs;
treating all empty queries as mutually related would therefore create
spurious positives.

$J$ is relative to a graph at a given point in time: a pair that is
high-overlap on the current version of DBpedia may not be so after a
subsequent release. We make this explicit by computing every label from a
cached snapshot of the answer sets; labels are thus reproducible without
endpoint access, and the construction can be repeated when the graph evolves.
Since the answers are collected from public endpoints, they reflect one
consistent graph only if the endpoint data does not change during collection.
All queries were executed against local mirrors of the two graphs
\hp{[TODO: name the DBpedia and Wikidata dump releases loaded in the mirrors
and the execution dates; do not name the hosting institution (anonymity).]}.

\paragraph{Learning objective.}
We seek an encoder $f$ that maps a query string to a unit vector $z = f(q)$
such that the cosine similarity $s(q_i,q_j) = z_i^{\top} z_j$ ranks high-overlap
candidates ahead of low-overlap candidates for $q_i$. This ranking objective
does not require cosine similarity to estimate $J$ numerically. We separately
assess whether a threshold $\theta$ can identify high-overlap pairs reliably
enough for approximate reuse. At inference time, $f$ has access to query
text only.

\subsection{\squareds{}: Answer Supervision from Query Logs}
\label{sec:data}

\squareds{} (\textbf{S}PARQL \textbf{Qu}ery \textbf{A}nswer \textbf{R}anking
\textbf{E}valuation) converts a query log into contrastive training instances.
Each training instance $(q_a,\mathcal{P},\mathcal{N})$ associates an anchor query $q_a$
with a set of positives $\mathcal{P}$ and a set of negatives $\mathcal{N}$.
We construct \squareds{} for DBpedia and Wikidata from the corresponding LSQ
logs~\cite{saleem2015lsq}; Table~\ref{tab:dataset_square} summarizes its
statistics.

\paragraph{Execution oracle.}
We sample \texttt{SELECT} queries from LSQ and keep those that return at
least one row. A first labeling pass compares the first 200 rows of each
answer and retains the anchors with at least one candidate positive,
together with their sampled negatives. We then execute every retained query
once more, without a row limit, and recompute all labels from the complete
answers. This selection favors queries that have a related query in the
log, which is the population our tasks address. On DBpedia, the retained
queries stem from 59{,}000 sampled LSQ queries, of which 16{,}450 returned
results \hp{[TODO: verify (the mining log reports 50{,}000 input queries) and
add the Wikidata counts]}. The
canonicalized answer sets are cached, and all labels are recomputed from this
cache rather than inherited from upstream files, so that every label remains
verifiable. Supervision is free of annotation but not of cost: it requires one
execution per logged query and one Jaccard computation per candidate pair.
On our hardware, executing all queries takes \hp{[TODO: --]} and labeling
all pairs \hp{[TODO: --]}; at inference time, encoding a query takes \hp{[TODO: --]}\,ms and a
nearest-neighbor search over the cached queries \hp{[TODO: --]}\,ms on
\hp{[TODO: hardware]}.

\paragraph{Component-disjoint splits.}
High overlap is symmetric, and the queries form a closed set. If an
anchor $a$ is held out while one of its positives $q$ remains in the training
set as an anchor, the pair $\{a,q\}$ is observed during training and evaluated
again at test time. We prevent this by construction: we build an undirected
graph whose edges are the positive pairs and assign each connected component
in its entirety to the training, validation or test split (80/10/10, fixed
seed). Consequently, queries that are linked directly or transitively never
appear in different splits. We audit the resulting splits at three levels: query strings, normalized
queries and answer sets. Normalization parses each query into its SPARQL
algebra, which expands prefixed IRIs, and renames variables in order of
appearance. No query string, positive pair or canonical answer set occurs in
two splits, and no normalized query is shared with the test split. The positive graph has 1{,}156 components on
DBpedia and 3{,}970 on Wikidata; most are small (median size 2), and the
largest contains 161 and 398 queries, respectively.
A component is a conservative grouping, not an equivalence class, and the
split does not rule out every shortcut. In particular, it does not prevent
the reuse of query \emph{templates}, i.e., queries that differ only in the
IRIs and literals in subject and object positions. Template reuse is
pervasive on DBpedia: its 4{,}727 queries follow only 121 templates, and 98\%
of the test anchors share a template with a training query, against 58\% on
Wikidata. Section~\ref{sec:rq2} therefore also reports results on test
anchors with unseen templates.

\paragraph{Instance construction.}
We construct instances within each split. Each anchor receives up to 20
positives, ranked by decreasing $J$, and up to 20 negatives. Among negatives,
\emph{nonzero-overlap} negatives that share at least one answer row with the
anchor ($0 < J \leq 0.1$) are preferred over fully disjoint ones ($J = 0$).
Nonzero-overlap negatives are taken in order of decreasing $J$, and the
remaining slots are filled with randomly sampled disjoint queries from the
same split. Nonzero-overlap negatives make up 38\% of the negatives on DBpedia
and 47\% on Wikidata. Nonzero answer overlap does not establish lexical or
structural difficulty, and some negatives are easy for a simpler reason: 25\%
(DBpedia) and 34\% (Wikidata) of them differ from their anchor in projection
arity, whereas almost all positives (over 99\%) share it
(Section~\ref{sec:rq2}). In contrast to negatives mined
with an embedding model~\cite{xiong2021ance}, which may be false negatives,
every explicit negative in \squareds{} is verified by the oracle for its
own anchor. This does not imply that all cross-triplet in-batch negatives are
low-overlap (Section~\ref{sec:objective}). Anchors without at
least one positive and one negative are discarded.

\begin{table}[t]
\centering
\caption{Statistics of \squareds{}. Positives satisfy $J \geq 0.8$ and
negatives $J \leq 0.1$; each anchor has at most 20 of each. Components are
those of the positive graph used for splitting.}
\label{tab:dataset_square}
\small
\setlength{\tabcolsep}{6pt}
\renewcommand{\arraystretch}{1.1}
\begin{tabular}{lrr}
\toprule
& \textbf{DBpedia} & \textbf{Wikidata} \\
\midrule
Queries & 4{,}727 & 10{,}749 \\
Anchors (train / valid / test) & 3{,}737 / 469 / 468 & 7{,}570 / 946 / 946 \\
Positive / negative pairs & 39{,}430 / 93{,}480 & 53{,}234 / 189{,}240 \\
Positive pairs with $J = 1$ & 79\% & 86\% \\
Components (largest) & 1{,}156 (161) & 3{,}970 (398) \\
Query templates & 121 & 5{,}509 \\
Test anchors with a training template & 98\% & 58\% \\
\bottomrule
\end{tabular}
\end{table}

\FloatBarrier
\subsection{Contrastive Fine-Tuning}
\label{sec:objective}

We fine-tune a pretrained query encoder with InfoNCE, using in-batch
candidates and execution-verified explicit
negatives~\cite{oord2018cpc,karpukhin2020dpr}.

A batch contains $B$ anchors. For anchor $q_i^a$, we sample one positive
$q_i^p$ from its stored positives and $m$ explicit negatives
$q_{i1}^n,\dots,q_{im}^n$ from its stored negatives. Every positive and
explicit negative in the batch serves as a candidate for every anchor, which
gives the candidate set
$\mathcal{C}=\{q_j^p\}_{j=1}^{B}\cup\{q_{jk}^n\}_{j\leq B,\,k\leq m}$ of size
$B(m+1)$. The encoder $f_\theta$ produces unit-normalized embeddings, and $s$
is their cosine similarity (Section~\ref{sec:problem}). For temperature
$\kappa>0$, the batch loss is
\begin{equation}
\mathcal{L}
=
-\frac{1}{B}\sum_{i=1}^{B}
\log
\frac{\exp\!\left(s(q_i^a,q_i^p)/\kappa\right)}
{\sum_{c\,\in\,\mathcal{C}\setminus\mathcal{M}_i}\exp\!\left(s(q_i^a,c)/\kappa\right)},
\label{eq:loss}
\end{equation}
where $\mathcal{M}_i$ contains the candidates other than $q_i^p$ that are
stored positives of $q_i^a$; masking them ensures that known positives never
act as negatives.

Execution verifies each anchor's own positive and explicit negatives.
Candidates contributed by other anchors are not verified for $q_i^a$, but
they can only be unlabeled positives if they lie in the positive component
of $q_i^a$ (Section~\ref{sec:data}). In a simulation of our sampler, this is
the case for at most 0.4\% (DBpedia) and 0.1\% (Wikidata) of these
candidates; intermediate overlap ($0.1<J<0.8$) remains possible.

\paragraph{Hard negatives.} In the first epoch, explicit negatives are
sampled at random. Before the second epoch, and every two epochs
thereafter, we encode the training queries with the current model and keep,
for each anchor, the 8 stored negatives with the highest cosine similarity;
three of the four explicit negatives are then drawn from this pool and one
at random. Unlike negatives mined from an index~\cite{xiong2021ance}, these
hard negatives are drawn from execution-verified negatives only, so they
cannot be false negatives.

The objective optimizes relative ranking: it favors each designated
positive over its batch competitors, but it does not require cosine
similarity to equal Jaccard overlap. We therefore introduce no
Jaccard-regression term and evaluate retrieval and threshold-based
approximate reuse separately (Section~\ref{sec:exp}).

To attribute the observed effects to the supervision, we use a shared,
minimal fine-tuning protocol for all encoders. We tokenize each raw query
string with the encoder's native tokenizer, pool its token representations
into a single vector, normalize the vector to unit length, and update all
encoder parameters.
Encoder models use mean pooling, except BGE and GTE, which use their
\texttt{[CLS]} vector; decoder models (Qwen3-Embedding, Qwen2.5-Coder,
Gemma-3) use the last token with left padding, and E5 receives its
\texttt{query:} prefix. All encoders share one configuration: $B=4$ anchors and $m=4$ explicit
negatives, so each anchor competes with 19 candidates, gradient accumulation
over 8 batches (32 anchors per update), one pass over all training anchors
per epoch, AdamW with
learning rate $10^{-5}$ and weight decay $0.01$, a cosine schedule with 5\%
warm-up, gradient clipping at $1.0$, $\kappa = 0.05$, 20 epochs, bf16 mixed
precision, and a maximum length of 512 tokens. Truncation can only affect long queries: no DBpedia query exceeds 1{,}145
characters, and 4.5\% of the Wikidata queries exceed 1{,}000 characters
\hp{[TODO: exact share above 512 tokens per tokenizer, scripts/audit/token\_lengths.py]}. We evaluate on the validation split after every
epoch and keep the checkpoint with the highest graded nDCG@10.
After training, we freeze the encoder and use query text alone to
rank cached queries.

% ===========================================================================
\section{Experiment}
\label{sec:exp}

We setup our experiments to answer the following questions: \textbf{i) RQ}$_1$.\textit{ Does fine-tuning on \squareds{} align embedding similarity with answer-set overlap?}, \textbf{ii) RQ}$_2$. D\textit{oes it improve the retrieval of high-overlap queries, compared with zero-shot encoders, non-neural baselines and fine-tuning on syntax-derived labels, and which relationships beyond syntax does it capture?} and \textbf{iii) RQ}$_3$. \textit{Can frozen embeddings decide when a cached result may be reused without query execution, and how often are its canonical answer sets identical?}

\subsection{Experimental Setup}
\label{sec:setup}

\paragraph{Pretrained encoders.}
We consider $10$ pretrained models from four families, selected to cover a
range of pretraining signals:
(i) \emph{text embedding models} trained contrastively on large text-pair
corpora, namely MPNet~\cite{song2020mpnet}, E5-base-v2~\cite{wang2022e5},
BGE-base~\cite{xiao2024bge}, GTE-base~\cite{li2023gte} and the LLM-based
Qwen3-Embedding~\cite{zhang2025qwen3emb};
(ii) a \emph{general-purpose language model} without embedding-specific
training, Gemma-3-270M~\cite{gemmateam2025gemma3};
(iii) \emph{code models}, namely CodeBERT~\cite{feng2020codebert},
GraphCodeBERT~\cite{guo2021graphcodebert} and
Qwen2.5-Coder~\cite{hui2024qwen25coder}; and
(iv) a \emph{SPARQL-pretrained model}, SPBERT~\cite{tran2021spbert}.

\paragraph{Non-neural baselines.}
We compare against TF-IDF cosine similarity and BM25~\cite{robertson2009bm25}
($k_1=1.2$, $b=0.75$) over SPARQL tokens, with variable names replaced by a
placeholder and statistics fitted on the training split, and against a
symbolic \emph{normalized match}: two queries match if their SPARQL algebra
coincides after prefix expansion and variable renaming. Normalized match
detects equality under this normalization, but does not identify arbitrary
rewrites or data-dependent overlap. As a structural comparator,
\emph{algebra similarity} ranks candidates by the weighted Jaccard
similarity of the multisets of algebra operators, triple patterns (with
variables abstracted) and parent--child operator pairs. A \emph{query-form}
baseline uses only shallow features: projection arity, number of triple
patterns, and the presence of each operator (\texttt{OPTIONAL},
\texttt{UNION}, \texttt{FILTER}, aggregates, solution modifiers), ranked by
Euclidean distance after standardization. It shows how much of the task can
be solved without any information about the data. A \emph{random ranking} of each anchor's
candidates gives the chance level of every metric.

\paragraph{Syntax-supervised control (planned comparison).}
To separate execution grounding from task-specific fine-tuning, the comparison
will use syntax-derived labels with the same encoder, objective and training
budget. Algebra similarity defines the labels, with per-anchor sampling used
to match positive and negative counts.
\hp{[TODO: run the matched-supervision comparison. Hold checkpoint,
objective, candidate counts, training steps and validation-selection rule
fixed. Prefer a shared candidate universe selected independently of answer
overlap. If both arms instead use execution-selected candidates, describe
the experiment as relabeling within an execution-selected pool; do not claim
that it isolates all effects of execution supervision. Fix the algebra
scorer and label-selection rule before test evaluation.]}

\paragraph{Protocol.}
Each encoder is evaluated in two settings: zero-shot, and after fine-tuning
on \squareds{} from the same checkpoint. Architecture, tokenizer, model
capacity and embedding dimensionality are held fixed, so that the supervision
is the only experimental variable. Fine-tuned results are reported as mean and
standard deviation over five training seeds. Seeds capture optimization
variance, but not the variance due to which query families were sampled.
Since anchors of one component share positives and failure modes,
uncertainty should also reflect the sampled query families.
\hp{[TODO: compute paired component-bootstrap 95\% confidence intervals
(10,000 resamples) for the principal FT--ZS and FT--strongest-baseline
differences. Resample whole components and preserve pairing between methods;
document seed aggregation and the number of test components. Report these
intervals separately from training-seed standard deviations.]}
All reported results use the designated held-out test split; its separation
is subject to the audit in Section~\ref{sec:data}.

\paragraph{Metrics.}
For RQ1, we report Spearman's rank correlation $\rho$ between the cosine
similarity $s$ and the Jaccard overlap $J$ over all test anchor--candidate
pairs. For RQ2, the candidates of each anchor are ranked by $s$, and a
candidate is relevant if $J \geq 0.8$. In the main setting, each anchor is ranked against
all other test queries (472 on DBpedia, 1{,}074 on Wikidata), and a candidate
counts as relevant only if it is one of the anchor's stored positives. This
undercounts relevant queries, since each anchor stores at most 20 positives;
the full-collection setting in Section~\ref{sec:rq2} uses complete labels
instead. \hp{[TODO: confirm that the reported neural scores used this
protocol (evaluation\_test.json, field candidate\_universe); the code on the
experiment branch does.]} We report Hits@$k$, MRR@$k$,
Precision@$k$, Recall@$k$ and F1@$k$ with binary relevance, and nDCG@$k$ with
$J$ as graded relevance, for $k \in \{1,3,5\}$, macro-averaged over anchors.
Under binary relevance, Hits@1, MRR@1 and Precision@1 coincide, and the
remaining metrics are strongly correlated, so they are not independent
evidence. The main text reports Hits@1, MRR@5 and nDCG@5; all metrics are given in
Appendix~\ref{app:full}. The caching metrics are defined in
Section~\ref{sec:caching}. All scores are reported with two decimal places.

\subsection{RQ1: Correlation with Answer Overlap}
\label{sec:rq1}

Table~\ref{tab:alignment} (Appendix~\ref{app:full}) reports the rank
correlation before and after fine-tuning. In the zero-shot setting, all encoders exhibit a moderate
correlation with answer overlap (mean $\rho = 0.70$ on DBpedia and $0.59$ on
Wikidata), reflecting the fact that queries sharing many answers often share
many tokens. Fine-tuning increases the mean correlation only marginally
($+0.02$ and $+0.03$, respectively), but substantially reduces the variation
across encoders, with the standard deviation decreasing from $0.06$ to $0.02$
on DBpedia and from $0.03$ to $0.02$ on Wikidata. Encoders with a low
zero-shot $\rho$ tend to gain more (Spearman $-0.87$ on both graphs), although
this is partly a statistical artifact, since the change is computed from the
zero-shot value itself. Gemma-3-270M, the weakest zero-shot encoder, exhibits the largest improvement
($+0.14$ and $+0.09$), whereas several well-aligned encoders decrease
slightly, for example GTE-base on DBpedia ($-0.05$).

This is consistent with the objective, which compares a designated positive
with the competitors in each training batch, rather than directly regressing
answer overlap. The pooled $\rho$ additionally compares similarities across
evaluation anchors. To test this
directly, we also compute $\rho$ within each anchor's candidate set and
average it over anchors: its measured change is \hp{[TODO: --]} to \hp{[TODO: --]} on
DBpedia and \hp{[TODO: --]} to \hp{[TODO: --]} on Wikidata.

RQ1 is thus only weak evidence that the learned similarity
\emph{estimates} $J$; what improves is ranking, which RQ2 measures directly.

\textit{Finding 1.} Answer supervision brings heterogeneous
pretrained representation spaces to a comparable level of answer awareness,
but changes the pooled correlation with answer overlap only marginally and
does not by itself induce a globally consistent similarity scale.

\subsection{RQ2: Retrieval of Related Queries}
\label{sec:rq2}

\input{Tables/compact/retrieval}

Fine-tuning improves retrieval for all ten encoders on both knowledge graphs
(Table~\ref{tab:retrieval}); the improvement holds for every metric and
cut-off in Appendix~\ref{app:full}, which, given the correlation between the
metrics, we regard as one consistent outcome rather than many. Seed standard
deviations are at most $0.02$ \hp{[TODO: insert the component-bootstrap 95\%
CI of the mean Hits@1 gain per graph]}. On DBpedia,
Hits@1 increases by $0.17$ to $0.29$ (mean $0.34 \rightarrow 0.55$) and nDCG@5
by $0.16$ to $0.26$; Qwen2.5-Coder more than doubles its Hits@1
($0.27 \rightarrow 0.56$). On Wikidata, the improvements are smaller but
consistent, with Hits@1 increasing by $0.04$ to $0.16$ (mean
$0.29 \rightarrow 0.39$). In absolute terms, the best fine-tuned encoder ranks
a high-overlap query first for $58\%$ of DBpedia and $41\%$ of Wikidata
anchors, against a chance level of $0.02$ and $0.01$, so the
task is far from solved. Lexical matching is a strong baseline: TF-IDF
reaches a Hits@1 of $0.58$ on DBpedia, on par with the best fine-tuned
encoder and above the fine-tuned mean ($0.55$), although fine-tuned encoders
rank the remaining positives better (nDCG@5 $0.55$ against $0.47$). On
Wikidata, TF-IDF ($0.36$) falls below every fine-tuned encoder
($0.38$--$0.41$). Fine-tuning thus brings zero-shot encoders (mean $0.34$
and $0.29$) to the level of lexical matching on DBpedia, where 98\% of the
test anchors share a template with training, and beyond it on Wikidata.
Algebra similarity and query-form features are weaker, and normalized match
recovers only 0.2\% (DBpedia) and 0.3\% (Wikidata) of the test positive pairs. Failure to match
under this normalization does not establish graph dependence, since other
semantics-preserving rewrites can yield distinct normalized strings.
The syntax-supervised comparison remains necessary to distinguish execution
supervision from task-specific fine-tuning more generally
\hp{[TODO: insert control scores and paired confidence intervals; state the
scope of attribution allowed by the candidate-selection design.]}



Three further observations follow. \emph{First, the improvement is inversely
related to zero-shot quality} (Figure~\ref{fig:gain} in Appendix~\ref{app:full}; as for RQ1, this correlation is partly an artifact of
computing the gain from the zero-shot value), and fine-tuning compresses the
differences between encoders: the range of Hits@1 narrows from
$0.27$--$0.41$ to $0.50$--$0.58$ on DBpedia, and from $0.23$--$0.37$ to
$0.38$--$0.41$ on Wikidata. \emph{Second, across the evaluated encoders and the shared protocol,
answer supervision substantially reduces the spread in retrieval
performance.} Since the encoders also differ in architecture, size, tokenizer
and training data, this does not isolate the effect of the pretraining domain.
Gemma-3-270M, a general-purpose language model without
embedding-specific training, attains a Hits@1 of $0.55$ on DBpedia, close to
the best encoder, while SPBERT, CodeBERT and GraphCodeBERT remain below all
five text embedding models after fine-tuning. \emph{Third, retrieval improves
even where pooled correlation does not.} GTE-base loses $0.05$ in $\rho$ on
DBpedia (Table~\ref{tab:alignment}) yet obtains one of the largest retrieval
gains ($+0.23$ Hits@1), which corroborates the within-anchor interpretation of
RQ1. Wikidata is more challenging for every encoder; a plausible explanation
is that its opaque identifiers (e.g., \texttt{wdt:P31}) provide text encoders
with little lexical evidence about the answers.

\paragraph{What does the supervision teach beyond syntax?}
This is the central question of the study: does execution supervision help
where lexical and structural evidence is misleading? The planned diagnostic separates three
categories of test pairs, each with inspected examples:
(i) \emph{validated rewrites}, pairs related by a semantics-preserving
rewrite, which we verify with \hp{[TODO: a set of algebra-level rewrite rules
/ manual inspection of a sample]}; they test robustness to alternative
formulations;
(ii) \emph{structural look-alikes}, pairs whose algebra expressions are
close (structural similarity $\geq \hp{[TODO: --]}$) but whose answers barely
overlap ($J \leq 0.1$); they test resistance to misleading syntactic
similarity. Queries (a) and (b) in Figure~\ref{fig:gap} resemble this
category, but with $J = 0.19$ they are an intermediate pair and do not belong
to it;
(iii) \emph{graph-dependent overlap}, high-overlap pairs ($J \geq 0.8$) that
are not related by any of the tested rewrites; their overlap is not explained
by the rewrites we test, although we cannot exclude that some are equivalent
on all graphs.
For each category, the comparison will use the zero-shot and fine-tuned encoders, the
syntax-supervised control and the structural comparator by the mean rank
percentile of the partner query among all other test queries.
\hp{[TODO: add a small table (3 categories $\times$ ZS / FT / syntax control /
structural comparator, both graphs) and one inspected example per category.
Report results whatever their direction. State category sizes, selection
rules and category overlap. Do not infer data dependence from normalized
match failure; optionally validate a subset with another graph snapshot or
a counterexample graph.]}

\paragraph{Shortcut audit.}
Some negatives may be easy for reasons unrelated to the data. Under
tuple-based comparison, answers with different projection arity can never
share a row, so an encoder could separate them by counting projected
variables. Likewise, two unrelated \texttt{COUNT} queries that return the same
integer have identical answer sets: such a pair is a valid positive under $J$,
but weak evidence of a shared information need. The shortcut audit should report Hits@1
(i) with negatives restricted to those with the anchor's projection arity and
query form, (ii) separately for aggregate and non-aggregate anchors, and
(iii) by answer-set size and number of triple patterns.
\hp{[TODO: run matched-arity/query-form retrieval and stratify by aggregate
status, singleton versus larger answer sets, and query size. Report stratum
sizes and FT--ZS/FT--baseline gains; state whether easy output-shape cues or
coincident aggregate values explain the gains.]}

\paragraph{Unseen templates.}
A split that is both component- and template-disjoint is not feasible: merging
queries along positive edges and shared templates yields one group with 61\%
of all DBpedia queries and one with 30\% of all Wikidata queries. We therefore
evaluate the existing models on the test anchors whose template does not occur
in training, 395 anchors on Wikidata and only 9 on DBpedia, so that this
analysis is informative for Wikidata only.
\hp{[TODO W2: report ZS/FT Hits@1 and nDCG@5 on the 395 unseen-template Wikidata
anchors next to the seen-template anchors (per-anchor scores already exist; no
retraining). State plainly if gains shrink.]}

\paragraph{Full-corpus retrieval.}
The main setting counts only stored positives as relevant. We therefore
also rank each test anchor against all other test queries with complete
labels, taken from the exhaustive pair labeling of the test split: positives
are rare (6.5\% of all test pairs on DBpedia and 0.8\% on Wikidata), and
intermediate pairs ($0.1 < J < 0.8$) count as non-relevant.
Table~\ref{tab:fullcollection} reports Hits@1, MRR and nDCG@10. With
complete labels, TF-IDF reaches a Hits@1 of $0.58$ on DBpedia and $0.38$ on
Wikidata. These collections contain hundreds to about a thousand queries and
do not establish large-scale indexing.
\input{Tables/compact/fullcollection}

\textit{Finding 2.} Answer supervision consistently improves the
retrieval of high-overlap queries across all model families, with the
largest gains for the weakest encoders. On DBpedia it reaches, but does
not exceed, a TF-IDF baseline; on Wikidata it moves beyond it.
\hp{[TODO: extend once the full-corpus, category and syntax-control results
are in.]}

\subsection{RQ3: Execution-Free Approximate Result Reuse}
\label{sec:caching}

\paragraph{Task and workload.}
This experiment diagnoses threshold-based decisions; it does not demonstrate
a deployed cache or correctness-preserving reuse. A semantic cache processes a stream of queries. For an incoming query $q$, it
retrieves the most similar cached query $c^* = \arg\max_{c \in C} s(q,c)$ and
reuses its result if $s(q,c^*) \geq \theta$ (a \emph{hit}); otherwise, $q$ is
executed and inserted under a FIFO policy with capacity $K = 128$. Answer
overlap is used only afterwards, to assess the decision. We replay all unique
test queries (473 for DBpedia, 1{,}075 for Wikidata) in five random orders and
label every query pair by $J$. The replay population is distinct from the
retrieval-anchor population (468 and 946, respectively, in
Table~\ref{tab:dataset_square}): retrieval requires a positive and a negative
for each anchor, whereas replay includes the full held-out query population.
The additional 5 and 129 queries are test queries without any high-overlap
partner in the test split. Positive pairs are considerably rarer on
Wikidata (0.8\% of all pairs, compared with 6.5\% on DBpedia). Replaying unique
queries isolates the reuse that string matching cannot find, but it removes
the repeated queries that make exact caching effective in real workloads; an
exact-match cache therefore has no hits here by design.
\hp{[TODO: add a replay of the original LSQ stream with repetitions (and
timestamps, if available), comparing an exact-match cache, a normalized-match
cache, and an exact-match cache backed by the learned cache for its misses.
Report the additional hits and their accuracy beyond exact caching.]}

\paragraph{Approximation contract.}
Reusing the answer of $c^*$ for $q$ is exact under our set-based answer
representation only if $J(q,c^*) = 1$ (Section~\ref{sec:problem}). We
evaluate an \emph{approximate} cache whose contract follows from the
positive threshold: if $J(q,c^*) \geq 0.8$, then the reused answer contains at
least $80\%$ of the rows of $\Ans(q,G)$, and at least $80\%$ of its rows belong
to $\Ans(q,G)$, since both $|A \cap B|/|A|$ and $|A \cap B|/|B|$ are bounded
below by $J$. Such reuse suits exploratory or recommendation settings that
tolerate small differences; applications that require exact SPARQL responses additionally need compatible
column mappings, multiplicities and ordering. Even $J = 1$ establishes only
identity of canonical answer sets, and no similarity threshold alone can
guarantee it. We
use reuse as a demanding test of the embeddings, not as a proposal for a
correctness-preserving cache.

\paragraph{Metrics.}
A hit is \emph{correct} if $J \geq 0.8$ (i.e., it satisfies the contract),
\emph{set-exact} if $J = 1$, \emph{false} if $J \leq 0.1$, and
\emph{intermediate} otherwise. With $N_q$ incoming queries, $N_h$ hits, of
which $N_c$ are correct, $N_e$ exact and $N_f$ false, we report the correct hit
rate $\mathrm{CHR} = N_c/N_q$ and the false hit rate $\mathrm{FHR} = N_f/N_q$.
Our primary reliability measure is the \emph{hit accuracy}
$\mathrm{Acc.} = N_c/N_h$, the share of all reused answers that satisfy the
contract. As a secondary diagnostic, the \emph{precision against false hits}
$\mathrm{Prec.} = N_c/(N_c+N_f)$ ignores intermediate hits and only measures
how well the cache avoids clearly wrong reuse; it is always at least as high
as Acc.\ and must not be read as the share of correct answers. Finally, the \emph{set-exact hit share} $N_e/N_h$ gives
the share of hits with identical canonical answer sets
\hp{[TODO: compute from the replay logs]}. An oracle that reuses a result
whenever a high-overlap query is cached attains a CHR of $0.72$ on
DBpedia and $0.37$ on Wikidata.

\input{Tables/compact/caching}

\paragraph{Aggregation and uncertainty.}
Table~\ref{tab:caching} reports a ratio of mean hit rates,
$\mathrm{Acc.}_{\mathrm{pool}} = \overline{\mathrm{CHR}} /
\overline{\mathrm{SHR}}$, corresponding to pooling hit counts across runs
with equal numbers of arrivals. This differs from a mean of per-run
accuracies and weights runs by their number of hits. The current accuracy
values are reconstructed from the reported two-decimal rates and should be
read as approximate summaries.
\hp{[TODO: recompute pooled accuracy from unrounded per-run hit counts;
report a confidence interval and the mean per-run accuracy, documenting
zero-hit runs and the treatment of seeds and stream orders.]}
Reuse is also sensitive to the run: on Wikidata, Qwen3-Embedding's secondary
precision is $0.45 \pm 0.31$ (Appendix~\ref{app:full}), so the largest mean
in the table does not establish a stable advantage.

\paragraph{Fixed threshold.}
We first use a single threshold $\theta = 0.8$ for all models, which exposes
differences in similarity scale (Table~\ref{tab:caching}). At this threshold,
zero-shot encoders reuse almost indiscriminately: their similarities are
concentrated in a narrow, high range, so nearly every query exceeds $\theta$
and the cache does not grow beyond its first entry, which explains why six
encoders obtain identical zero-shot scores. This shows a scale mismatch, not
that zero-shot encoders cannot serve as caches at a better threshold. After
fine-tuning, the hit accuracy on DBpedia lies between $0.59$ and $0.77$: between
23\% and 41\% of the reused answers violate the contract. Clearly wrong reuse
becomes rare (false hit rate between $0.01$ and $0.11$, precision against false
hits between $0.77$ and $0.96$), so most violations are intermediate hits. On
Wikidata, $\theta = 0.8$ remains too permissive: no encoder reaches a hit
accuracy above $0.28$.

\paragraph{Matched-accuracy comparison.}
A fixed threshold favors encoders whose scale happens to suit it. The fairer
comparison asks how many correct reuse decisions each method enables at the
same accuracy requirement. For each method, including the zero-shot encoders
and the lexical and structural baselines, we select $\theta$ on the validation
stream by replaying it independently at every candidate threshold and
maximizing validation CHR subject to validation hit accuracy of at least $0.9$.
Ties are broken by higher validation accuracy, then by the larger threshold.
Thresholds with no hits are ineligible. If no threshold is feasible, the
method abstains from reuse; its hit rate is zero and hit accuracy is undefined.
We report test CHR and achieved test accuracy at the selected threshold
(Appendix~\ref{app:calibration}); the validation constraint does not guarantee
that test accuracy reaches $0.9$.
\hp{[TODO: insert test CHR and achieved test accuracy per method and graph
after validation selection; report failures to meet 0.9 on test explicitly.]}

\textit{Finding 3.} Ranking quality and similarity calibration are distinct
problems: the same fine-tuned encoder that retrieves well can yield unreliable
reuse decisions under a fixed threshold. Even on DBpedia, a substantial share
of reused answers misses the overlap criterion, so learned similarity alone
does not make reuse reliable, let alone exact.

\subsection{Limitations}
\label{sec:limitations}

\emph{Graph-relative labels.} The relatedness $J$ is defined on a single
snapshot of each graph; the construction must be repeated when the graph
changes substantially, and transfer across graph versions is not evaluated.
\emph{In-batch candidates.} Known positives are masked and at most 0.4\% of
the in-batch candidates can be unlabeled positives, but candidates with
intermediate overlap are still treated as negatives.
\emph{Empty answers.} Queries with empty answers are excluded, so the encoder
is not trained to recognize equivalent queries that both return nothing.
\emph{Template overlap.} Most DBpedia test anchors (98\%) share a template
with a training query, so the DBpedia results mainly measure generalization
to new constants within known query shapes; only Wikidata supports an
unseen-template analysis (Section~\ref{sec:rq2}).
\emph{Approximate reuse.} A positive pair shares at least $80\%$ of its
answers, not all of them; our cache evaluation therefore concerns approximate
reuse. Exact SPARQL response reuse needs compatible output semantics and
verification beyond an embedding threshold. \emph{Offline caching evaluation.} The replay uses
unique queries only and measures decision quality, not realized latency or the
gain over an exact cache on a workload with repetitions. \emph{Scope.} We study
the DBpedia and Wikidata logs of LSQ; cross-graph transfer is limited by
disjoint vocabularies and is left to future work.

% ===========================================================================
\section{Conclusion}
\label{sec:conclusion}

We introduced answer-supervised representation learning for SPARQL: query
embeddings that rank queries by the overlap of their answers on a given graph,
learned
from supervision obtained by executing a query log once. Embedding-based
ranking then requires only query text; a reuse system still executes cache
misses. The \squareds{} resource provides this
supervision for DBpedia and Wikidata with component-disjoint splits. Across
ten pretrained encoders, answer-supervised fine-tuning consistently improves
retrieval of high-overlap queries over zero-shot checkpoints and
substantially reduces the spread between encoders. This establishes a
benefit of the complete fine-tuning procedure in the evaluated setting,
not yet a causal advantage of execution labels over syntax supervision.
\hp{[TODO: update this conclusion after full-collection baseline,
matched-supervision and generalization results are available. State measured
advantages and limitations; do not claim overlap estimation, equivalence
recognition or learning beyond syntax from the current results alone.]} Our results further show that ranking quality and
similarity calibration are distinct problems: improved retrieval does not
ensure reliable result reuse at a fixed threshold, and exact SPARQL response
reuse requires verification beyond embedding similarity.
Promising directions include objectives
that learn from the graded Jaccard
values in \squareds{} rather than from binary labels, generalization across
graph versions and query templates, and query logs with more complex queries.

\subsubsection*{Resource Availability.}
\label{sec:availability}
\squareds{} (training instances, cached answer sets and splits), the trained checkpoints
and all code are intended to be provided at
\url{https://anonymous.4open.science/r/SQuARE-ESWC27} for review, and
published under CC~BY~4.0 (data) and MIT (code) with a persistent DOI upon
acceptance. The queries derive from LSQ~\cite{saleem2015lsq}, whose
attribution requirements are propagated. Since the answer sets are cached,
the supervision remains usable even if the source endpoints change or become
unavailable.
\hp{[TODO: create the anonymized mirror \texttt{SQuARE-ESWC27} and remove
identifying terms from the code (see \texttt{CODE\_PAPER\_AUDIT.md}). Verify
public reviewer access and every released dependency, then update the
availability statements to describe the actual release.]}

% ===========================================================================
\bibliographystyle{splncs04}
\bibliography{references}

% ===========================================================================
\clearpage
\appendix
\renewcommand{\theHsection}{appendix.\Alph{section}}
\section{Full Results}
\label{app:full}

Table~\ref{tab:alignment} gives the RQ1 correlations per encoder and
Figure~\ref{fig:gain} the retrieval gain against zero-shot quality.
Tables~\ref{tab:dbpedia_results}--\ref{tab:wikidata_binary} report all
retrieval metrics for $k \in \{1,3,5\}$, and
Tables~\ref{tab:app-caching-db} and~\ref{tab:app-caching-wd} report the
complete caching results, including the total hit rate (SHR) and the
fraction of available reuse opportunities exploited by the cache
(Opp.~Rec.). The intermediate hit rate is $\mathrm{SHR} - \mathrm{CHR} -
\mathrm{FHR}$, and the hit accuracy of Table~\ref{tab:caching} is
$\mathrm{CHR}/\mathrm{SHR}$.

\input{Tables/compact/alignment}
\begin{figure}[htbp]
\centering
\begin{tikzpicture}
\begin{axis}[
  width=0.62\textwidth, height=4.1cm,
  xlabel={Zero-shot Hits@1}, ylabel={Gain in Hits@1},
  xlabel style={font=\scriptsize}, ylabel style={font=\scriptsize},
  tick label style={font=\scriptsize},
  xmin=0.2, xmax=0.43, ymin=0, ymax=0.32,
  xticklabel style={/pgf/number format/.cd, fixed, fixed zerofill, precision=2},
  yticklabel style={/pgf/number format/.cd, fixed, fixed zerofill, precision=2},
  legend style={font=\scriptsize, at={(0.98,0.97)}, anchor=north east, draw=black!30},
  grid=major, grid style={black!10},
]
\addplot[only marks, mark=*, mark size=1.8pt, color=blue!70!black]
  table[x=zs, y=gain] {Tables/compact/gain_db.dat};
\addlegendentry{DBpedia}
\addplot[only marks, mark=triangle*, mark size=2.3pt, color=orange!85!black]
  table[x=zs, y=gain] {Tables/compact/gain_wd.dat};
\addlegendentry{Wikidata}
\end{axis}
\end{tikzpicture}
\caption{Gain from fine-tuning against zero-shot Hits@1, one point per
encoder (Spearman $-0.87$ on DBpedia, $-0.99$ on Wikidata). The
correlation is descriptive: the gain is computed from the zero-shot value.}
\label{fig:gain}
\end{figure}

\input{Tables/appendix/retrieval_full}
\input{Tables/appendix/caching_full}

\clearpage
\section{Calibrated Thresholds for Result Reuse}
\label{app:calibration}

Retrieval only requires a correct ordering of candidates, whereas a semantic
cache must decide whether a single similarity value is high enough to reuse a
stored result. Since Equation~\ref{eq:loss} does not train similarity to estimate $J$,
a threshold that suits one encoder or graph need not suit another
(Section~\ref{sec:caching}). We therefore select the threshold $\theta$ on
held-out data, with two procedures that use the frozen embeddings and the
validation split only.

\emph{Validation-optimal threshold.} For each candidate $\theta$, we reset
and replay the validation workload under the same cache policy. We choose the
threshold maximizing validation CHR subject to validation hit accuracy of at
least $0.9$, using the tie-breaking and abstention rules of
Section~\ref{sec:caching}. A new replay is necessary because each threshold
changes the miss decisions and hence the subsequent cache contents.
The grid is $\theta\in\{0.50,0.60,0.70,0.75,0.80,0.85,0.90,0.95\}$, and
hit counts are pooled over five validation stream orders. Test accuracy is measured after selection and may be below $0.9$.

\emph{Conformal risk control (CRC).} As a second procedure, we adopt
conformal risk control~\cite{angelopoulos2024crc,vovk2005alrw}. For each of
$n$ validation arrivals $i$, we replay the validation stream with the same
cache policy and let $c_i$ be the most similar query in the cache \emph{as it
stands when $q_i$ arrives}, with similarity $s_i$. We define the loss
$L_i(\theta) = \mathbf{1}[s_i \geq \theta \wedge J(q_i,c_i) < 0.8]$, which
indicates a hit that violates the contract, intermediate hits included. With
the empirical risk $\hat{R}_n(\theta) = \frac{1}{n}\sum_i L_i(\theta)$, we set
\begin{equation}
\hat{\theta} = \inf\Big\{\theta : \tfrac{n}{n+1}\hat{R}_n(\theta) + \tfrac{1}{n+1} \leq \alpha\Big\}.
\label{eq:crc}
\end{equation}
Under the assumptions of conformal risk control, this loss would control
the joint probability $P(\text{hit} \wedge \text{violation})$, averaged over
arrivals, rather than the conditional violation probability among hits,
$P(\text{violation} \mid \text{hit}) = 1 - \mathrm{Acc.}$ If the joint
probability were bounded by $\alpha$ and $P(\text{hit}) > 0$, the conditional
probability would be bounded only by
$\min\{1, \alpha/P(\text{hit})\}$. A cache that reuses on $5\%$ of arrivals
and is always wrong still meets a per-arrival bound of $0.05$.
The theorem additionally requires monotonic losses and exchangeability.
These assumptions are not established for our replay: cache contents vary
with $\theta$, and arrivals are evaluated using evolving shared state. We
therefore use Equation~\ref{eq:crc} only as a \emph{CRC-inspired heuristic}
targeting per-arrival violations, and claim no formal risk bound for this
adaptive replay. The validation-optimal procedure directly targets hit
accuracy and is the main calibration comparison.
Neither procedure requires retraining or access to
answers at inference time. Besides single operating points, we report
risk--coverage curves, i.e., $1 - \mathrm{Acc.}$ against the hit rate as
$\theta$ varies \hp{[TODO: add one plot per graph for three encoders]}.

\begin{table}[t]
\centering
\caption{Calibrated thresholds for result reuse. Fixed: $\theta = 0.8$;
Val-opt: highest correct hit rate on validation with hit accuracy $\geq 0.9$;
CRC: CRC-inspired heuristic targeting per-arrival contract violations
with $\alpha = 0.05$ (Eq.~\ref{eq:crc}); no formal replay guarantee is claimed.
Fine-tuned encoders; per training seed, the five stream orders are
averaged, and we report the mean over five seeds.}
\label{tab:calibration}
\small
\setlength{\tabcolsep}{4pt}
\renewcommand{\arraystretch}{1.1}
\begin{adjustbox}{max width=\textwidth}
\begin{tabular}{ll cccc p{0pt} cccc}
\toprule
& & \multicolumn{4}{c}{\textbf{DBpedia}} & & \multicolumn{4}{c}{\textbf{Wikidata}}\\
\cmidrule(lr){3-6}\cmidrule(lr){8-11}
Encoder & Threshold & Prec.\,$\uparrow$ & Acc.\,$\uparrow$ & CHR\,$\uparrow$ & FHR\,$\downarrow$ & & Prec.\,$\uparrow$ & Acc.\,$\uparrow$ & CHR\,$\uparrow$ & FHR\,$\downarrow$\\
\midrule
E5-base-v2 & Fixed & 0.96 & 0.77 & 0.20 & 0.01 & & 0.23 & 0.22 & 0.14 & 0.47 \\
 & Val-opt & \hp{--} & \hp{--} & \hp{--} & \hp{--} & & \hp{--} & \hp{--} & \hp{--} & \hp{--} \\
 & CRC & \hp{--} & \hp{--} & \hp{--} & \hp{--} & & \hp{--} & \hp{--} & \hp{--} & \hp{--} \\
\midrule
Qwen3-Embedding & Fixed & 0.89 & 0.67 & 0.28 & 0.05 & & 0.45 & 0.28 & 0.08 & 0.21 \\
 & Val-opt & \hp{--} & \hp{--} & \hp{--} & \hp{--} & & \hp{--} & \hp{--} & \hp{--} & \hp{--} \\
 & CRC & \hp{--} & \hp{--} & \hp{--} & \hp{--} & & \hp{--} & \hp{--} & \hp{--} & \hp{--} \\
\midrule
CodeBERT & Fixed & 0.79 & 0.61 & 0.37 & 0.10 & & 0.23 & 0.22 & 0.14 & 0.48 \\
 & Val-opt & \hp{--} & \hp{--} & \hp{--} & \hp{--} & & \hp{--} & \hp{--} & \hp{--} & \hp{--} \\
 & CRC & \hp{--} & \hp{--} & \hp{--} & \hp{--} & & \hp{--} & \hp{--} & \hp{--} & \hp{--} \\
\bottomrule
\end{tabular}
\end{adjustbox}
\end{table}

Table~\ref{tab:calibration} reports the results.
\hp{[TODO: report measured CHR, achieved hit accuracy, and the per-arrival
violation rate, including intermediate hits. Compare calibration methods
without presuming that either meets its target on test.]}

\end{document}
